\subsection{半局部环的 Hausdorff 完备化}

命题 17。设 A 为交换环，\((\mathfrak{m}_{\lambda})_{\lambda \in L}\) 是 A 的一个理想族，各理想均不等于 A，并且当 \(\mathfrak{m}_{\lambda}\) 与 \(\mathfrak{m}_{\mu}\) 满足 \(\lambda \neq \mu\) 时二者互素。对于任意一个族 \(s = (s(\lambda))_{\lambda \in L}\)，其元素为 \(\geqslant 0\) 的整数且具有有限支撑，置 \(\mathfrak{a}_{s} = \bigcap_{\lambda \in L} \mathfrak{m}_{\lambda}^{s(\lambda)}\)（它等于所有 \(\mathfrak{m}_{\lambda}^{s(\lambda)}\) 的乘积，其中的 \(\lambda\) 满足 \(s(\lambda) \neq 0\)；参见~第二章，第 1 节，第 2 号，命题 3 和 5）；\(\mathfrak{a}_{s}\) 构成关于与 A 的环结构相容的拓扑 \(\mathcal{T}\) 的 0 的邻域基本系；设 \(\hat{\mathbf{A}}\) 为 A 关于此拓扑的 Hausdorff 完备化。另一方面，对于所有 \(\lambda \in L\)，令 \(\mathbf{A}_{\lambda}\) 为赋予 \(\mathfrak{m}_{\lambda}\)-进拓扑的环 A，并令 \(\hat{\mathbf{A}}_{\lambda}\) 为其 Hausdorff 完备化。若 \(u: \mathbf{A} \to \prod_{\lambda \in L} \mathbf{A}_{\lambda}\) 表示对角同态，则 \(u\) 连续，且相应的同态 \(\hat{u}\)：

\[
\hat {\mathbf {A}} \rightarrow \left(\prod_ {\lambda \in \mathbf {L}} \mathbf {A} _ {\lambda}\right) ^ {\wedge} = \prod_ {\lambda \in \mathbf {L}} \hat {\mathbf {A}} _ {\lambda}
\]

（一般拓扑学，第三章，第 6 节，第 5 号以及第二章，第 3 节，第 9 号，命题 18 的推论 2）是一个拓扑环同构。

第一个断言由一般拓扑学第三章，第 6 节，第 3 号，例 3 得出。置 \(\mathbf{B} = \prod_{\lambda \in \mathbf{L}} \mathbf{A}_{\lambda}\)；由于 \(\mathbf{A}\) 上的拓扑比每个 \(m_{\lambda}\)-进拓扑都细，映射 \(\mathrm{pr}_{\lambda} \circ u\) 连续，因而 \(u\) 连续。此外，\(u(a_s)\) 是 \(\Delta\)（\(\mathbf{B}\) 的对角线）与开集 \(\bigcap_{\lambda \in \mathbf{L}} \mathrm{pr}_{\lambda}^{-1}(m_{\lambda}^{s(\lambda)})\)（\(\mathbf{B}\) 中）的交；由此可知，\(u\) 是从加法群 \(\mathbf{A}\) 到 \(\mathbf{B}\) 的严格态射，其像为 \(\Delta\)。现在 \(\Delta\) 在 \(\mathbf{B}\) 中稠密。事实上，设 \(b = (a_{\lambda})_{\lambda \in \mathbf{L}}\) 是 \(\mathbf{B}\) 的一个元素；\(b\) 在 \(\mathbf{B}\) 中的每个邻域都包含形如 \(b + \mathbf{V}\) 的集合，其中 \(\mathbf{V} = \bigcap_{\lambda \in \mathbf{L}} \mathrm{pr}_{\lambda}^{-1}(m_{\lambda}^{s(\lambda)})\)，而 \(s = (s(\lambda))_{\lambda \in \mathbf{L}}\) 是一个具有有限支撑的、由整数 \(\geqslant 0\) 构成的族。由于 \(m_{\lambda}^{s(\lambda)}\) 两两互素（第二章，第 1 节，第 2 号，命题 3），存在 \(x \in \mathbf{A}\)，使得 \(x \equiv a_{\lambda}\)（模 \(m_{\lambda}^{s(\lambda)}\)）对所有 \(\lambda\) 都成立（同上，命题 5），从而 \((b + \mathbf{V}) \cap \Delta \neq \varnothing\)。于是群 \(\mathbf{B}/\Delta\) 的 Hausdorff 完备化为 \(\{0\}\)；将第 12 号的引理 2 应用于正合列 \(0 \to \mathbf{A} \xrightarrow{u} \mathbf{B}, \mathbf{A} \xrightarrow{u} \mathbf{B} \to \mathbf{B}/\Delta\)，可见 \(\hat{u}\) 是将 \(\hat{\mathbf{A}}\) 同构到 \(\hat{\mathbf{B}}\) 的同构。

推论。设 A 为主理想整环，P 为 A 的极端元代表系（代数，第七章，第 1 节，第 3 号）。A 上使 A 的 \(\neq0\) 理想构成 0 的邻域基本系、且与 A 的环结构相容的拓扑是 Hausdorff 的；A 关于此拓扑的完备化典范同构于 A 关于 \(\pi\)-进拓扑的完备化的乘积，其中 \(\pi\) 遍历 P。

主理想 \((\pi)\)（其中 \(\pi \in P\)）是极大的且彼此不同，因而互素；我们已经看到（第 12 号，例 3），\(\pi\)-进拓扑是 Hausdorff 的，所以命题 17 的叙述中定义的拓扑也是 Hausdorff 的，因为它比每个 \(\pi\)-进拓扑都细。

将命题 17 的推论应用于 A = Z 时，我们以 \(\hat{Z}\) 表示 Z 关于使 Z 的所有 \(\neq0\) 理想构成 0 的邻域基本系的拓扑的完备化；这个环同构于 p-进整数环的乘积 \(\prod_{p\in P}Z_{p}\)，其中 P 是素数集。

\textbf{注}\par

\begin{enumerate}
\def\labelenumi{(\arabic{enumi})}
\item
  显然，在命题 17 的条件下，拓扑 \(\mathcal{T}\) 是各 \(m_{\lambda}\)-进拓扑在 \(\mathbf{A}\) 上的上确界。
\item
  \(\mathfrak{a}\) 是 \(\prod_{\lambda \in L} \hat{\mathbf{A}}_{\lambda}\) 的每个闭理想，都等于其投影 \(\mathfrak{a}_{\lambda} = \operatorname{pr}_{\lambda}(\mathfrak{a})\) 的乘积，而这些投影是 \(\hat{\mathbf{A}}_{\lambda}\) 中的闭理想；因为 \(\hat{\mathbf{A}}_{\lambda}\) 典范地识别为闭理想 \(\mathbf{A}_{\lambda}'\)，它是 \(\prod_{\lambda} \hat{\mathbf{A}}_{\lambda}\) 的闭理想，并且 \(\mathfrak{a}_{\lambda}\) 识别为 \(\mathfrak{a} \cap \mathbf{A}_{\lambda}'\)（代数，第一章，第 8 节，第 10 号，命题 6），\(\mathfrak{a}_{\lambda}\) 的和在乘积 \(\prod_{\lambda} \mathfrak{a}_{\lambda}\) 中稠密（一般拓扑学，第三章，第 2 节，第 9 号，命题 25），而后者在 \(\prod_{\lambda} \hat{\mathbf{A}}_{\lambda}\) 中闭，由此得到所述断言。
\end{enumerate}

命题 18。设 A 为交换环，\((\mathfrak{m}_i)_{1 \leqslant i \leqslant q}\) 是 A 的有限个互不相同的极大理想组成的族，r 是乘积理想 \(\mathfrak{m}_1\mathfrak{m}_2 \ldots \mathfrak{m}_q = \mathfrak{m}_1 \cap \mathfrak{m}_2 \cap \cdots \cap \mathfrak{m}_q\)，S 是乘性子集 \(\bigcap_{i=1}^{q} (\mathbf{A} - \mathfrak{m}_i)\)。赋予 A r-进拓扑，赋予环 \(\mathbf{B} = \mathbf{S}^{-1}\mathbf{A}\) 以 rB-进拓扑，并赋予每个局部环 \(\mathbf{A}_{\mathfrak{m}_i}\) 以 \((\mathfrak{m}_i\mathbf{A}_{\mathfrak{m}_i})\)-进拓扑。令 \(u: \mathbf{A} \to \mathbf{B}\)、\(v_i: \mathbf{B} \to \mathbf{A}_{\mathfrak{m}_i}\) 为典范同态（第二章，第 2 节，第 1 号，命题 2 的推论 2），令 \(v\) 为同态 \((v_i): \mathbf{B} \to \prod_{i=1}^{q} \mathbf{A}_{\mathfrak{m}_i}\)。同态 \(u\) 和 \(v\) 连续，相应的同态 \(\hat{u}: \hat{\mathbf{A}} \to \hat{\mathbf{B}}\) 和 \(\hat{v}: \hat{\mathbf{B}} \to \prod_{i=1}^{q} (\mathbf{A}_{\mathfrak{m}_i})^{\wedge}\) 是拓扑环同构。

\(m_{i} \cap S = \varnothing\)（\(1 \leqslant i \leqslant q\)），所以 B 的理想 \(m_{i}' = m_{i}B\) 是极大的（第二章，第 2 节，第 5 号，命题 11），且

\[
\mathfrak {r B} = \mathfrak {m} _ {1} ^ {\prime} \cap \mathfrak {m} _ {2} ^ {\prime} \cap \dots \cap \mathfrak {m} _ {q} ^ {\prime}
\]

(第二章，第 2 节，第 4 号）；最后，\(\mathbf{B}_{\mathfrak{m}_i'} = \mathbf{A}_{\mathfrak{m}_i}\) 在一个典范同构下成立（第二章，第 2 节，第 5 号，命题 11）。由于 \(\bar{u}^{-1}(\mathfrak{rB}) = \mathfrak{r}\)，且 \(\bar{v}_i^{-1}(m_i\mathbf{A}_{\mathfrak{m}_i}) \supset \mathfrak{rB}\)，\(u\) 和 \(v\) 连续。因此只需证明：

\[
w = v \circ u: \mathrm{A} \rightarrow \prod_ {i = 1} ^ {q} \mathrm{A} _ {\mathfrak {m} _ {i}},
\]

\(\hat{w}\) 是 \(\hat{\mathbf{A}}\) 到 \(\prod_{i=1}^{q} \hat{\mathbf{A}}_{m_i}\) 上的同构，因为将此结果应用于 B 和 \(m_i'\) 将表明 \(\hat{v}\) 是同构，因而 \(\hat{u}\) 也是同构。注意，每个 \(m_i\) 的幂的乘积都包含 r 的某个幂，所以 r-进拓扑是各 \(m_i\)-进拓扑的上确界；此外，若 \(A_i\) 表示

\[
\mathfrak {m} _ {i} \text {-adic}
\]

\[
\phi : \mathrm{A} \rightarrow \prod_ {i = 1} ^ {q} \mathrm{A} _ {i}
\]

\(\hat{\phi}:\hat{\mathbf{A}}\to \prod_{i = 1}^{q}\hat{\mathbf{A}}_i\) 是一个同构（命题 17）。于是问题归结为证明：若 \(u_{i}:\mathbf{A}_{i}\rightarrow \mathbf{A}_{m_{i}}\) 是典范映射，则 \(\hat{u}_i\colon \hat{\mathbf{A}}_i\to \hat{\mathbf{A}}_{m_i}\) 是同构。现在，对所有 \(n\)，映射

\[
u _ {i, n} \colon \mathrm{A} / \mathfrak {m} _ {i} ^ {n} \rightarrow \mathrm{A} _ {\mathfrak {m} _ {i}} / \mathfrak {m} _ {i} ^ {n} \mathrm{A} _ {\mathfrak {m} _ {i}}
\]

由 \(u_{i}\) 取商得到的映射是同构（第二章，第 3 节，第 3 号，命题 9）；我们的断言由以下事实得出：\(\hat{\mathbf{A}}_i\)（分别为 \(\hat{\mathbf{A}}_{\mathfrak{m}_i}\)）是离散环 \(\mathbf{A} / \mathfrak{m}_i^n\)（分别为 \(\mathbf{A}_{\mathfrak{m}_i} / \mathfrak{m}_i^n \mathbf{A}_{\mathfrak{m}_i}\)）的逆极限（参见~第 6 号）。

注（3）。我们看到，一个整环 A 的 Hausdorff 完备化 \(\hat{A}\) 可能含有非零的零因子。

命题 19。设 A 为交换环，m 为 A 的极大理想。A 关于 m-进拓扑的 Hausdorff 完备化 \(\hat{\mathbf{A}}\) 是一个局部环，其极大理想为 \(\hat{\mathfrak{m}}\)。

若 \(\mathfrak{a} = \bigcap_{k\geqslant 1}\mathfrak{m}^k\)，则 \(\hat{\mathbf{A}}\) 是 Hausdorff 环 \(\mathbf{A} / \mathfrak{a}\)（与 \(\mathbf{A}\) 相伴）的完备化；又由于 \(\mathfrak{m} / \mathfrak{a}\) 是 \(\mathbf{A} / \mathfrak{a}\) 中的极大理想，我们可设 \(\mathbf{A}\) 关于 \(\mathfrak{m}\)-进拓扑是 Hausdorff 的。由于 \(\mathbf{A} / \mathfrak{m}\) 与 \(\hat{\mathbf{A}} /\hat{\mathfrak{m}}\) 是同构环（第 12 号，公式（21）），\(\hat{\mathfrak{m}}\) 是 \(\hat{\mathbf{A}}\) 的极大理想。由于 \(\hat{\mathbf{A}}\) 上的拓扑由滤过 \((\mathfrak{m}^n)^{\wedge}\) 定义（第 12 号），命题将由下述引理推出：

引理 3。设 A 为完备的 Hausdorff 拓扑环，且其中存在一个由 A 的加法子群组成的 0 的邻域基本系 \(\mathfrak{S}\)。

\begin{enumerate}
\def\labelenumi{(\roman{enumi})}
\item
  对于所有 \(x \in \mathbf{A}\) 满足 \(\lim_{n \to \infty} x^n = 0\)，\(1 - x\) 在 \(\mathbf{A}\) 中可逆，且其逆等于 \(\sum_{n=0}^{\infty} x^n\)。
\item
  设 \(\mathfrak{a}\) 是 A 的一个双边理想，满足 \(\lim x^n = 0\) 对所有 \(x \in \mathfrak{a}\) 都成立。A 的元素 y 可逆的充要条件是其模 a 的类在 A/a 中可逆；特别地，a 包含于 A 的 Jacobson 根中。
\end{enumerate}

(i) 由于

\[
(1 - x) (1 + x + \dots + x ^ {n}) = (1 + x + \dots + x ^ {n}) (1 - x) = 1 - x ^ {n + 1},
\]

问题归结为证明通项为 \(x^{n}\) 的级数在 A 中收敛；现在，根据假设，对于 A 中 0 的每个邻域 \(V \in \mathfrak{S}\)，存在整数 p \textgreater{} 0，使得 \(x^{n} \in V\) 对所有 \(n \geqslant p\) 都成立。于是

\[
x ^ {p} + x ^ {p + 1} + \dots + x ^ {q} \in \mathbf {V}
\]

对所有 \(q \geqslant p\) 都成立，而我们的断言随后由柯西判别法得出（一般拓扑学，第三章，第 5 节，第 2 号，定理 1）。

\begin{enumerate}
\def\labelenumi{(\roman{enumi})}
\setcounter{enumi}{1}
\tightlist
\item
  假设存在 \(y' \in A\)，使得 \(yy' \equiv 1 \pmod{a}\) 且 \(y'y \equiv 1 \pmod{a}\)。由 (i)，关于 \(a\) 的假设蕴含 \(yy'\) 和 \(y'y\) 在 \(A\) 中可逆，因而 \(y\) 在 \(A\) 中可逆。特别地，对每个 \(x \in a\)，\(1 - x\) 在 \(A\) 中可逆；又由于 \(a\) 是 \(A\) 的双边理想，它包含于 \(A\) 的 Jacobson 根中（代数，第八章，第 6 节，第 3 号，定理 1）。
\end{enumerate}

既已证明此引理，只需将其应用于拓扑环 \(\hat{\mathbf{A}}\) 及理想 \(\hat{\mathfrak{m}}\)，因为对于所有 \(x \in \hat{\mathfrak{m}}\)，都有 \(x^n \in (\hat{\mathfrak{m}})^n \subset (\mathfrak{m}^n)^\wedge\)，所以序列 \((x^n)\) 趋于 0。

若取 A = Z，则 Z 的每个极大理想都形如 pZ，其中 p 为素数。于是 p-进数环 \(Z_{p}\) 是一个局部环，其极大理想为 \(pZ_{p}\)（命题 16 的推论 2），其剩余域同构于 \(Z/pZ = F_{p}\)；而 \(\mathbf{Z}_{(p)}\) 赋予 \(pZ_{(p)}\)-进拓扑后，可识别为 \(Z_{p}\) 的一个包含 Z 的拓扑子环。

推论。设 A 为半局部环（第二章，第 3 节，第 5 号），\(\mathfrak{m}_i\) 为其互不相同的极大理想（\(1 \leqslant i \leqslant q\)），且

\[
\mathfrak {r} = \mathfrak {m} _ {1} \cap \mathfrak {m} _ {2} \cap \dots \cap \mathfrak {m} _ {q}
\]

它的 Jacobson 根。Hausdorff 完备化 \(\hat{\mathbf{A}}\) 是 \(\mathbf{A}\) 关于 \(\mathfrak{r}\)-进拓扑的完备化，是一个半局部环，典范同构于乘积 \(\prod_{i=1}^{q} \hat{\mathbf{A}}_{\mathfrak{m}_i}\)，其中 \(\hat{\mathbf{A}}_{\mathfrak{m}_i}\) 是局部环 \(\mathbf{A}_{\mathfrak{m}_i}\) 关于 \((\mathfrak{m}_i\mathbf{A}_{\mathfrak{m}_i})\)-进拓扑的 Hausdorff 完备化环。

